\documentclass[12pt]{article}
\usepackage{amsmath,amssymb,amsthm,physics,geometry}
\geometry{margin=1in}

\title{Spin in Nelson-Type Stochastic Mechanics:\\
From Pauli Theory to Relativistic Spin Bundles}
\author{}
\date{}

\begin{document}
\maketitle

\begin{abstract}
We present a systematic treatment of spin within Nelson-type stochastic
mechanics. Starting from the scalar diffusion model reproducing the
Schr\"odinger equation, we construct three progressively deeper
extensions: (1) a stochastic variational derivation of the Pauli equation,
(2) an internal diffusion model on $SU(2)$ generating half-integer
representations geometrically, and (3) a principal-bundle formulation
leading naturally toward the Dirac equation in the relativistic setting.
\end{abstract}

\tableofcontents

% ==========================================================
\section{Scalar Nelson Mechanics}

Let $X_t \in \mathbb{R}^3$ satisfy the It\^o SDE
\begin{equation}
dX_t = b(X_t,t) dt + \sqrt{\frac{\hbar}{m}}\, dW_t.
\end{equation}

Define forward/backward drifts $b_\pm$ and
\begin{equation}
v = \frac{b_+ + b_-}{2}, \qquad
u = \frac{b_+ - b_-}{2}.
\end{equation}

The density $\rho$ satisfies
\begin{equation}
\partial_t \rho + \nabla\cdot(\rho v)=0.
\end{equation}

Consistency of forward/backward diffusion yields
\begin{equation}
u = \frac{\hbar}{2m}\nabla \ln \rho.
\end{equation}

Imposing mean Newton equation
\begin{equation}
m\left(\partial_t v + (v\cdot\nabla)v - (u\cdot\nabla)u
- \frac{\hbar}{2m}\Delta u\right)
= -\nabla V,
\end{equation}
and setting $v=\frac{1}{m}\nabla S$, we obtain
\begin{align}
\partial_t \rho + \nabla\cdot\left(\rho \frac{\nabla S}{m}\right)&=0,\\
\partial_t S + \frac{(\nabla S)^2}{2m}
+ V - \frac{\hbar^2}{2m}\frac{\Delta\sqrt{\rho}}{\sqrt{\rho}}&=0.
\end{align}

Defining $\psi=\sqrt{\rho}e^{iS/\hbar}$ yields the Schr\"odinger equation.

This construction produces scalar (spin-0) theory only.

% ==========================================================
\section{Stochastic Derivation of the Pauli Equation}

Let $\Psi(x,t)\in\mathbb{C}^2$.

Define:
\begin{align}
\rho &= \Psi^\dagger\Psi,\\
s^i &= \Psi^\dagger\sigma^i\Psi.
\end{align}

The Pauli current is
\begin{equation}
j =
\frac{\hbar}{m}\Im(\Psi^\dagger\nabla\Psi)
- \frac{e}{m}A\rho
+ \frac{\hbar}{2m}\nabla\times s.
\end{equation}

We impose the stochastic continuity equation
\begin{equation}
\partial_t \rho + \nabla\cdot(\rho v)=0,
\end{equation}
with
\begin{equation}
v=\frac{j}{\rho}.
\end{equation}

Thus
\begin{equation}
v=\frac{1}{m}(\nabla S - eA)
+ \frac{\hbar}{2m}\frac{\nabla\times s}{\rho}.
\end{equation}

\subsection{Stochastic Variational Principle}

Consider
\begin{equation}
\mathcal{A}
=
E\int
\left[
\frac{m}{2}(v^2+u^2)
- e\phi
+ \frac{e}{2m}s\cdot B
\right] dt.
\end{equation}

Using
\[
(\sigma\cdot\Pi)^2
=
\Pi^2 - e\hbar\sigma\cdot B,
\]
variation yields
\begin{equation}
i\hbar\partial_t\Psi
=
\left[
\frac{1}{2m}(\sigma\cdot(-i\hbar\nabla-eA))^2
+ e\phi
\right]\Psi.
\end{equation}

Hence spin enters through a rotational correction to drift and
magnetic coupling in the stochastic action.

% ==========================================================
\section{Diffusion on $SU(2)$}

Let configuration space be
\[
Q=\mathbb{R}^3\times SU(2).
\]

Let $(X_t,g_t)$ satisfy coupled diffusions:
\begin{align}
dX_t &= b\,dt + \sqrt{\frac{\hbar}{m}}\,dW_t,\\
dg_t &= g_t \circ d\Xi_t,
\end{align}
where $\Xi_t$ is Brownian motion in $\mathfrak{su}(2)$.

Generator:
\begin{equation}
\mathcal{L}
=
\frac{\hbar}{2m}\Delta_{\mathbb{R}^3}
+
\kappa \Delta_{SU(2)}.
\end{equation}

Since
\[
\Delta_{SU(2)} = J^2,
\]
its eigenfunctions are irreducible representations labelled by $j$.

Restricting to $j=\frac12$ produces two-component spinors.

Thus spin emerges as internal Brownian motion on a compact Lie group.

% ==========================================================
\section{Principal Bundle Formulation}

Let $P(\mathbb{R}^3,SU(2))$ be the principal bundle.

Lift diffusion horizontally using connection $\omega$:
\[
\tilde{X}_t \in P.
\]

Magnetic interaction arises from curvature of a $U(1)$ connection:
\[
F=dA.
\]

Spin precession equals parallel transport in $SU(2)$ bundle.

Hence:
\[
\text{Spin} = \text{Holonomy of lifted stochastic flow}.
\]

% ==========================================================
\section{Hydrodynamic / Vorticity Interpretation}

Allow rotational component in velocity field:
\[
\omega=\nabla\times v.
\]

Spin density may be identified as
\[
s \propto \rho\,\omega.
\]

Then Pauli magnetic term represents coupling between vorticity
and electromagnetic curvature.

Spin becomes intrinsic circulation of diffusion flow.

% ==========================================================
\section{Relativistic Lift Toward Dirac}

Let diffusion occur in spacetime:
\[
dX^\mu = b^\mu d\tau + \sqrt{\frac{\hbar}{m}} dW^\mu.
\]

Lift to spin bundle over Minkowski space.
Introduce spin connection $\omega_\mu$ and covariant derivative
\[
D_\mu = \partial_\mu + \omega_\mu.
\]

Imposing relativistic stochastic Newton law yields
\[
(i\gamma^\mu D_\mu - m)\Psi=0.
\]

Spin now arises from representation theory of the Lorentz group
and the double cover $SL(2,\mathbb{C})$.

% ==========================================================
\section{Structural Synthesis}

Spin in stochastic mechanics may arise via:

\begin{enumerate}
\item Modified drift with rotational correction (Pauli level),
\item Internal diffusion on $SU(2)$,
\item Principal-bundle holonomy,
\item Intrinsic vorticity of stochastic flow,
\item Lorentz spin-bundle structure in relativistic theory.
\end{enumerate}

The deepest geometric statement is:

\[
\boxed{
\text{Spin = geometry of lifted stochastic diffusion on a principal bundle.}
}
\]

\end{document}
